\begin{lemma}
\label{lemma-geometrically-unibranch}
\begin{reference}
\cite[Lemma 2.2]{Etale-coverings} and
\cite[Chapter IV Proposition 18.8.15]{EGA4}
\end{reference}
Let $A$ be a local ring. Let $A^{sh}$ be a strict henselization of $A$.
The following are equivalent
\begin{enumerate}
\item $A$ is geometrically unibranch, and
\item $A^{sh}$ has a unique minimal prime.
\end{enumerate}
\end{lemma}

\begin{proof}
This follows from Lemma \ref{lemma-geometric-branches}
but we will also give a direct proof; this direct
proof is almost exactly the same as the direct proof of
Lemma \ref{lemma-unibranch}.
Denote $\mathfrak m$ the maximal ideal of the ring $A$.
Denote $\kappa$, $\kappa^{sh}$ the residue field of $A$, $A^{sh}$.

\medskip\noindent
Assume (2). Let $\mathfrak p^{sh}$ be the unique minimal prime of
$A^{sh}$. The flatness of $A \to A^{sh}$ implies that
$\mathfrak p = A \cap \mathfrak p^{sh}$ is the unique minimal
prime of $A$ (by going down, see
Algebra, Lemma \ref{algebra-lemma-flat-going-down}).
Also, since $A^{sh}/\mathfrak pA^{sh} = (A/\mathfrak p)^{sh}$ (see
Algebra, Lemma \ref{algebra-lemma-quotient-strict-henselization})
is reduced by Lemma \ref{lemma-henselization-reduced}
we see that $\mathfrak p^{sh} = \mathfrak pA^{sh}$.
Let $A'$ be the integral closure of $A/\mathfrak p$ in its fraction
field. We have to show that $A'$ is local and that its residue
field is purely inseparable over $\kappa$.
Since $A \to A'$ is integral, every maximal ideal of $A'$ lies over
$\mathfrak m$ (by going up for integral ring maps, see
Algebra, Lemma \ref{algebra-lemma-integral-going-up}).
If $A'$ is not local, then we can find distinct maximal ideals
$\mathfrak m_1$, $\mathfrak m_2$. Choosing elements $f_1, f_2 \in A'$
with $f_i \in \mathfrak m_i, f_i \not \in \mathfrak m_{3 - i}$ we find
a finite subalgebra $B = A[f_1, f_2] \subset A'$ with distinct maximal
ideals $B \cap \mathfrak m_i$, $i = 1, 2$. If $A'$ is local with maximal
ideal $\mathfrak m'$, but $A/\mathfrak m \subset A'/\mathfrak m'$
is not purely inseparable, then we can find $f \in A'$ whose image in
$A'/\mathfrak m'$ generates a finite, not purely inseparable extension
of $A/\mathfrak m$ and we find a finite local subalgebra $B = A[f] \subset A'$
whose residue field is not a purely inseparable extension of $A/\mathfrak m$.
Note that the inclusions
$$
A/\mathfrak p \subset B \subset \kappa(\mathfrak p)
$$
give, on tensoring with the flat ring map $A \to A^{sh}$ the inclusions
$$
A^{sh}/\mathfrak p^{sh} \subset
B \otimes_A A^{sh} \subset
\kappa(\mathfrak p) \otimes_A A^{sh} \subset
\kappa(\mathfrak p^{sh})
$$
the last inclusion because
$\kappa(\mathfrak p) \otimes_A A^{sh} =
\kappa(\mathfrak p) \otimes_{A/\mathfrak p} A^{sh}/\mathfrak p^{sh}$
is a localization of the domain $A^{sh}/\mathfrak p^{sh}$.
Note that $B \otimes_A \kappa^{sh}$ has at least two maximal ideals
because $B/\mathfrak mB$ either has two maximal ideals or one whose
residue field is not purely inseparable over $\kappa$, and because
$\kappa^{sh}$ is separably algebraically closed. Hence, as
$A^{sh}$ is strictly henselian we see that
$B \otimes_A A^{sh}$ is a product of $\geq 2$ local rings, see
Algebra, Lemma \ref{algebra-lemma-mop-up-strictly-henselian}.
But we've just seen that $B \otimes_A A^{sh}$ is a subring of a domain
and we get a contradiction.

\medskip\noindent
Assume (1). Let $\mathfrak p \subset A$ be the unique minimal
prime and let $A'$ be the integral closure of $A/\mathfrak p$
in its fraction field. Let $A \to B$ be a local map of local rings which is a
localization of an \'etale $A$-algebra. In particular $\mathfrak m_B$
is the unique prime containing $\mathfrak m_AB$. Then $B' = A' \otimes_A B$
is integral over $B$ and the assumption that $A \to A'$ is local
with purely inseparable residue field extension implies that $B'$
is local (Algebra, Lemma \ref{algebra-lemma-local-tensor-with-integral}).
On the other hand, $A' \to B'$ is the localization
of an \'etale ring map, hence $B'$ is normal, see
Algebra, Lemma \ref{algebra-lemma-normal-goes-up}.
Thus $B'$ is a (local) normal domain. Finally, we have
$$
B/\mathfrak pB \subset B \otimes_A \kappa(\mathfrak p)
= B' \otimes_{A'} (\text{fraction field of }A') \subset
\text{fraction field of }B'
$$
Hence $B/\mathfrak pB$ is a domain, which implies that $B$ has a unique
minimal prime (since by flatness of $A \to B$ these all have to lie
over $\mathfrak p$). Since $A^{sh}$ is a filtered colimit of
the local rings $B$ it follows that $A^{sh}$ has a unique minimal prime.
Namely, if $fg = 0$ in $A^{sh}$ for some non-nilpotent elements
$f, g$, then we can find a $B$ as above containing both $f$ and $g$
which leads to a contradiction.
\end{proof}